\subsection{Parallel Gates}

Now that we have a way to represent the state of multiple qubits, we can also represent the operations (gates) that act on them. In particular, if we have two qubits $\ket{v_{A}}$ and $\ket{v_{B}}$, and two gates $A$ and $B$, what happens when we want to apply $A$ to $\ket{v_{A}}$ and $B$ to $\ket{v_{B}}$ simultaneously?

\highspace
The answer is given by the \textbf{tensor product of quantum gates}:
\begin{equation}
    \left(A \otimes B\right)\left(\ket{v_{A}} \otimes \ket{v_{B}}\right) = \left(A\ket{v_{A}}\right) \otimes \left(B\ket{v_{B}}\right)
\end{equation}
This means that the combined operation of applying $A$ to $\ket{v_{A}}$ and $B$ to $\ket{v_{B}}$ is represented by the tensor product of the two gates, $A \otimes B$. So, \textbf{applying gates in parallel} corresponds to the tensor product of the individual gates.

\begin{figure}[!htp]
    \centering

    % Parallel single-qubit gates
    \begin{quantikz}
        \lstick{$\ket{v_A}$} & \gate{A} & \push{A\ket{v_A}} \rstick[2]{\hspace{3em}$(A\ket{v_A}) \otimes (B\ket{v_B})$} \\
        \lstick{$\ket{v_B}$} & \gate{B} & \push{B\ket{v_B}}
    \end{quantikz}

    \begin{equation*}
        \downarrow \hspace{3em} \text{Tensor Product} \hspace{3em} \downarrow
    \end{equation*}

    % Global tensor-product gate
    \begin{quantikz}
        \lstick{$\ket{v_A}$} & \gate[2]{A \otimes B} & \rstick[2]{$\left(A \otimes B\right)\ket{v_{A}v_{B}} = \left(A\ket{v_A}\right) \otimes \left(B\ket{v_B}\right)$} \\
        \lstick{$\ket{v_B}$} &                       &
    \end{quantikz}
    

    \caption{Parallel gates as a tensor product operation.}
\end{figure}

\highspace
In simple terms, instead of applying $A$ and $B$ separately, we build one big matrix $A \otimes B$ that acts on the \textbf{whole system}. Note that a state is represented as a column vector (e.g., size $4$ for two qubits), and a gate is represented as a square matrix (e.g., size $4 \times 4$ for two qubits). So \emph{parallel gates} are simply one big transformation of the global state of the system.

\highspace
\begin{flushleft}
    \textcolor{Red2}{\faIcon{question-circle} \textbf{What if we want to apply a gate to only one qubit?}}
\end{flushleft}
Suppose we want to \textbf{apply a gate} $A$ to the \textbf{first qubit} $\ket{v_{A}}$ while \textbf{leaving the second qubit} $\ket{v_{B}}$ \textbf{unchanged}. We can represent this as:
\begin{equation}
    \left(A \otimes I\right)\left(\ket{v_{A}} \otimes \ket{v_{B}}\right) = \left(A\ket{v_{A}}\right) \otimes \ket{v_{B}}
\end{equation}
Here, $I$ is the \textbf{identity gate} (\autopageref{sec:identity-gate}) that acts on the second qubit, meaning it \textbf{leaves it unchanged}. So, to apply a gate to just one qubit in a multi-qubit system, we use the tensor product of that gate with the identity operator on the other qubits.

\highspace
\begin{examplebox}[: $H \otimes X$]
    Let's consider the gates $H$ (Hadamard, \autopageref{subsubsec:hadamard-gate-h}):
    \begin{equation*}
        H = \frac{1}{\sqrt{2}}
        \begin{pmatrix}
            1 & 1 \\
            1 & -1
        \end{pmatrix}
    \end{equation*}
    And $X$ (Pauli-X, \autopageref{sec:pauli-x-not-gate}):
    \begin{equation*}
        X = \begin{pmatrix}
            0 & 1 \\
            1 & 0
        \end{pmatrix}
    \end{equation*}
    The tensor product $H \otimes X$ is calculated as follows:
    \begin{align*}
        H \otimes X =& \frac{1}{\sqrt{2}} \left(
            \begin{pmatrix}
                1 & 1 \\
                1 & -1
            \end{pmatrix}
            \otimes
            \begin{pmatrix}
                0 & 1 \\
                1 & 0
            \end{pmatrix}
        \right) \\
        =& \frac{1}{\sqrt{2}} \begin{pmatrix}
            1 \cdot X & 1 \cdot X \\
            1 \cdot X & -1 \cdot X
        \end{pmatrix} \\
        =& \frac{1}{\sqrt{2}} \begin{pmatrix}
            1 \cdot \begin{pmatrix}
                0 & 1 \\
                1 & 0
            \end{pmatrix} & 1 \cdot \begin{pmatrix}
                0 & 1 \\
                1 & 0
            \end{pmatrix} \\%[1em]
            1 \cdot \begin{pmatrix}
                0 & 1 \\
                1 & 0
            \end{pmatrix} & -1 \cdot \begin{pmatrix}
                0 & 1 \\
                1 & 0
            \end{pmatrix}
        \end{pmatrix} \\
        =& \frac{1}{\sqrt{2}} \begin{pmatrix}
            0 & 1 & 0 & 1 \\
            1 & 0 & 1 & 0 \\
            0 & 1 & 0 & -1 \\
            1 & 0 & -1 & 0
        \end{pmatrix}
    \end{align*}
    Now, if we apply $H \otimes X$ to a simple two-qubit state, say $\ket{00}$:
    \begin{align*}
        (H \otimes X)\ket{00} =& \frac{1}{\sqrt{2}} \begin{pmatrix}
            0 & 1 & 0 & 1 \\
            1 & 0 & 1 & 0 \\
            0 & 1 & 0 & -1 \\
            1 & 0 & -1 & 0
        \end{pmatrix} \begin{pmatrix}
            1 \\
            0 \\
            0 \\
            0
        \end{pmatrix} \\
        =& \frac{1}{\sqrt{2}} \begin{pmatrix}
            0 \cdot 1 + 1 \cdot 0 + 0 \cdot 0 + 1 \cdot 0 \\
            1 \cdot 1 + 0 \cdot 0 + 1 \cdot 0 + 0 \cdot 0 \\
            0 \cdot 1 + 1 \cdot 0 + 0 \cdot 0 + -1 \cdot 0 \\
            1 \cdot 1 + 0 \cdot 0 + -1 \cdot 0 + 0 \cdot 0
        \end{pmatrix} \\
        =& \frac{1}{\sqrt{2}} \begin{pmatrix}
            0 \\
            1 \\
            0 \\
            1
        \end{pmatrix} \\
        =& \frac{1}{\sqrt{2}} \begin{pmatrix}
            0 \\
            1 \\
            0 \\
            0
        \end{pmatrix} + \frac{1}{\sqrt{2}} \begin{pmatrix}
            0 \\
            0 \\
            0 \\
            1
        \end{pmatrix} \\
        =& \frac{1}{\sqrt{2}} \ket{01} + \frac{1}{\sqrt{2}} \ket{11}
    \end{align*}
    This can be visually represented as parallel gates acting on the two qubits:
    \begin{center}
        \begin{quantikz}
            \lstick{$\ket{0}$} & \gate{H} & \push{H\ket{0}} \rstick[2]{\hspace{2.15em}$\left(H\ket{0}\right) \otimes \left(X\ket{0}\right) = \ket{+} \otimes \ket{1}$} \\
            \lstick{$\ket{0}$} & \gate{X} & \push{X\ket{0}}
        \end{quantikz}
    \end{center}
    Or as a single global gate:
    \begin{center}
        \begin{quantikz}
            \lstick{$\ket{0}$} & \gate[2]{H \otimes X} & \rstick[2]{$\left(H \otimes X\right)\ket{00} = \ket{+1}$} \\
            \lstick{$\ket{0}$} &                       &
        \end{quantikz}
    \end{center}
    Where $\ket{+1}$ is $\ket{+} \otimes \ket{1}$. Note that applying the Hadamard gate to the qubit $\ket{0}$ creates the superposition state $\ket{+} = \frac{1}{\sqrt{2}}(\ket{0} + \ket{1})$, while applying the $X$ gate to $\ket{0}$ flips it to $\ket{1}$.
\end{examplebox}

\begin{examplebox}[: $I \otimes X$]
    Now, let's consider applying the identity gate $I$ to the first qubit and the $X$ gate to the second qubit. The tensor product is:
    \begin{align*}
        I \otimes X =& \begin{pmatrix}
            1 & 0 \\
            0 & 1
        \end{pmatrix} \otimes \begin{pmatrix}
            0 & 1 \\
            1 & 0
        \end{pmatrix} \\
        =&
        \begin{pmatrix}
            1 \cdot \begin{pmatrix}
                0 & 1 \\
                1 & 0
            \end{pmatrix} & 0 \cdot \begin{pmatrix}
                0 & 1 \\
                1 & 0
            \end{pmatrix} \\
            0 \cdot \begin{pmatrix}
                0 & 1 \\
                1 & 0
            \end{pmatrix} & 1 \cdot \begin{pmatrix}
                0 & 1 \\
                1 & 0
            \end{pmatrix}
        \end{pmatrix} \\
        =& \begin{pmatrix}
            0 & 1 & 0 & 0 \\
            1 & 0 & 0 & 0 \\
            0 & 0 & 0 & 1 \\
            0 & 0 & 1 & 0
        \end{pmatrix}
    \end{align*}
    Applying this to the state $\ket{00}$ gives:
    \begin{align*}
        (I \otimes X)\ket{00} =& \begin{pmatrix}
            0 & 1 & 0 & 0 \\
            1 & 0 & 0 & 0 \\
            0 & 0 & 0 & 1 \\
            0 & 0 & 1 & 0
        \end{pmatrix} \begin{pmatrix}
            1 \\
            0 \\
            0 \\
            0
        \end{pmatrix} \\
        =& \begin{pmatrix}
            0 \cdot 1 + 1 \cdot 0 + 0 \cdot 0 + 0 \cdot 0 \\
            1 \cdot 1 + 0 \cdot 0 + 0 \cdot 0 + 0 \cdot 0 \\
            0 \cdot 1 + 0 \cdot 0 + 0 \cdot 0 + 1 \cdot 0 \\
            0 \cdot 1 + 0 \cdot 0 + 1 \cdot 0 + 0 \cdot 0
        \end{pmatrix} \\
        =& \begin{pmatrix}
            0 \\
            1 \\
            0 \\
            0
        \end{pmatrix} \\
        =& \ket{01}
    \end{align*}
    This means that applying $I$ to the first qubit leaves it unchanged as $\ket{0}$, while applying $X$ to the second qubit flips it from $\ket{0}$ to $\ket{1}$, resulting in the combined state $\ket{01}$. Visually, this can be represented as:
    \begin{center}
        \begin{quantikz}
            \lstick{$\ket{0}$} & \gate{I} & \push{I\ket{0}} \rstick[2]{\hspace{2.15em}$\ket{0} \otimes (X\ket{0}) = \ket{0} \otimes \ket{1}$} \\
            \lstick{$\ket{0}$} & \gate{X} & \push{X\ket{0}}
        \end{quantikz}
    \end{center}
    Or as a single global gate:
    \begin{center}
        \begin{quantikz}
            \lstick{$\ket{0}$} & \gate[2]{I \otimes X} & \rstick[2]{$(I \otimes X)\ket{00} = \ket{01}$} \\
            \lstick{$\ket{0}$} &                       &
        \end{quantikz}
    \end{center}
\end{examplebox}